\documentclass[12pt,reqno]{amsart}
\usepackage{/Users/johnmyers/GoogleDrive/tex/handout,/Users/johnmyers/GoogleDrive/tex/customcommands,polynom}
\usepackage{caption}

\hdr{MAT 240}{\S16.5 {\small Integrals in cylindrical and spherical coordinates, part 2}}

\begin{document}

\textbf{Suggested problems:} 1-41 odd

\bigskip
\hrule

\bigskip
\prob Using the diagram below, write down equations that relate the spherical coordinates $(\rho,\theta,\phi)$ of a point to its cartesian (rectangular) coordinates $(x,y,z)$.

\medskip
\includegraphics[scale=1.25]{1.pdf}

\textcolor{red}{From spherical to cartesian coordinates we have:
	\[x = \rho  \cos{(\theta)}\sin{(\phi)}, \quad y = \rho\sin{(\theta)}\sin{(\phi)}, \quad z = \rho\cos{(\phi)}.
	\]
We also have $\rho^2 = x^2+y^2+z^2$.}

\medskip
\prob 

\begin{enumerate}
\item Plot the point with spherical coordinates $(2,\pi/4,\pi/3)$ and find its cartesian coordinates.

\bigskip
\textcolor{red}{See class notes for plot. Cartesian coordinates $= (\sqrt{6}/2,\sqrt{6}/2,1)$.}
\bigskip

\item Find the spherical coordinates of the point with cartesian coordinates $(0,2\sqrt{3},-2)$.

\bigskip
\textcolor{red}{We have $(\rho,\theta,\phi) = (4,\pi/2,2\pi/3)$.}
\bigskip
\end{enumerate}

\bigskip
\prob

\begin{enumerate}
\item Describe the equation whose equation in spherical coordinates is $\rho=1$. \textcolor{red}{Unit sphere centered at the origin.}
\item Give an equation in spherical coordinates for the cone $z = \sqrt{x^2+y^2}$. \textcolor{red}{$\phi = \pi/4$.}
\item Give an equation in spherical coordinates for the plane $z=10$. \textcolor{red}{$\rho \cos{(\phi)} = 10$.}
\end{enumerate}

\newpage
\prob Using the diagram below, express the volume element $dV$ in spherical coordinates.

\medskip

\includegraphics[scale=1.25]{pic2.pdf}

\bigskip
\textcolor{red}{$dV = \rho^2 \sin{(\phi)} d\phi d\theta d\rho$}
\bigskip

\prob For each object below, set up a triple integral of a function $f$ over the object.

\medskip
\begin{enumerate}
\item \textcolor{white}{blank} \\
\includegraphics[scale=1.25]{pic6.pdf}
\textcolor{red}{
	\[\int_0^{\pi}\int_0^\pi \int_2^3 f(\rho,\theta,\phi)\rho^2\sin \phi \ d\rho d \phi d\theta
	\]}
\item A piece of a sphere, angle at the center is $\pi/3$. \\

\includegraphics[scale=1.25]{2.pdf} \\
\textcolor{red}{
	\[\int_0^{\pi/6} \int_0^{2\pi} \int_0^3 f(\rho,\theta,\phi) \rho^2 \sin \phi \ d\rho d \theta d\phi 
	\]}	
\item \textcolor{white}{blank} \\
\includegraphics[scale=1.25]{pic7.pdf}
\textcolor{red}{
	\[\int_{\pi/4}^{\pi/2} \int_0^\pi \int_0^{2/\sin{(\phi)}} f(\rho,\theta,\phi)\rho^2\sin \phi \ d \rho d \theta d\phi
	\]}
\end{enumerate}

\bigskip
\prob Evaluate
	\[\int_W \sin{(\phi)} \ dV
	\]
where $W$ is the region $0 \leq \theta \leq 2 \pi$, $0 \leq \phi \leq \pi/4$, and $1 \leq \rho \leq 2$.

\bigskip
\textcolor{red}{
\begin{align*}
\int_0^{2\pi} \int_0^{\pi/4} \int_1^2 \sin^2{(\phi)}\rho^2 \ d\rho d\phi d \theta &= \frac{7}{3} \int_0^{2\pi} \int_0^{\pi/4} \sin^2{(\phi)} \ d\phi d\theta \\
&= \frac{7}{3} \int_0^{2\pi} \int_0^{\pi/4} \left( \frac{1}{2} - \frac{1}{2} \cos{(2\phi)} \right) \ d\phi d\theta \\
&= \frac{7}{3}\left( \frac{\pi}{8} - \frac{1}{4}\right) \int_0^{2\pi}  d\theta \\
&= \frac{14\pi}{3}\left( \frac{\pi}{8} - \frac{1}{4}\right).
\end{align*}}
\bigskip

\prob Compute the volume of the solid pictured below.

\medskip
\includegraphics[scale=1.25]{3.pdf}

\medskip
\textcolor{red}{
\begin{align*}
\int_0^{\pi/2} \int_{\pi/2}^{2\pi} \int_1^2 \sin{(\phi)} \rho^2 \ d\rho d\theta d\phi &= \frac{7}{3} \int_0^{\pi/2} \int_{\pi/2}^{2\pi} \sin{(\phi)} \ d\theta d \phi \\
&= \frac{7\pi}{2} \int_0^{\pi/2} \sin{(\phi)} \ d\phi \\
&= \frac{7\pi}{2}
\end{align*}}
\bigskip

\prob Evaluate
	\[\int_W (x^2+y^2) \ dV,
	\]
where $W$ lies between the spheres $x^2+y^2+z^2=4$ and $x^2+y^2+z^2=9$.

\bigskip
\textcolor{red}{
\begin{align*}
\int_W (x^2+y^2) \ dV &= \int_0^\pi \int_0^{2\pi} \int_2^3 (\rho^2 \sin^2{(\phi)}) \rho^2 \sin{(\phi)} \ d\rho d\theta d\phi \\
&= \frac{211}{5}\int_0^\pi \int_0^{2\pi} \sin^3{(\phi)} \ d\theta d\phi \\
&=  \frac{422\pi}{5} \int_0^\pi \sin^3{(\phi)} \ d\phi \\
&= \frac{422\pi}{5}\int_0^\pi(\sin{(\phi)} - \sin{(\phi)} \cos^2{(\phi)}) \ d\phi \\
&= \frac{1688\pi}{15}
\end{align*}}

\end{document}